% !TeX root = ../papers/nc-p3.tex
\section{二元商の全例外値、商の一意性と同型}
\label{sec:binary-apolar-rank-test}
\begin{theorem}\label{thm:binary-apolar-rank-test}
前節の $([T],[\bar Y])$ に対し、$I_6=\{f\in P_6:T(fP_4)=0\}$ とおく。
$\bar Y$ は自然に $(I_6\otimes P_2)^*$ の非零元をスカラーまで定める。
この二行三列行列が階数2であることと、構成した窓が有限Frobenius検査を通ることは同値である。
このとき反対階数は常に8である。階数1では三次交叉が既に過大となる。
従って有限点グラフを持つ非正則Frobenius窓の指定の向きでのパラメータは、
階数1のSegre多様体と前定理の $\pi(S_2)$ を除けば、例外なく全て許される。

指定の向きで二元商の核直線は一意で、二元座標は同時の $\operatorname{PGL}_2$ まで一意である。
有限グラフのこの向きの二組の表示が添字保持同型である必要十分条件は、
$([T],[\bar Y])$ が同じ $\operatorname{PGL}_2$ 軌道に属することである。
従ってこの向きの非正則有限境界は既約で、基底変更を除いた次元は12である。
\end{theorem}
\begin{proof}[差分式]
$P_d$ を一変数 $x$ の次数 $d$ 以下の空間として書く。モニック六次式 $f\in I_6$ を選べる。
実際次数5以下の係数に関する五条件の行は可逆Hankel行列 $G$ の初めの五行である。
$e\in P_5$ を $G(e,x^j)=\delta_{j5}$ で定め、斉次化ではこれを次数6にも見る。
$f,e$ は $I_6$ の基底で、共通零点を持たない。
実際 $A=k[x]/f$ で $T$ の次数10以下の積評価は剰余の評価だけで決まる
（$T(fP_4)=0$）。従って $G(e,b)$ は次数5以下の剰余 $b$ の最高次係数である。
もし $a$ が共通根なら $b=f/(x-a)$ に対し $eb=0$ が $A$ で成立し、その最高次係数1と矛盾する。
無限遠では $f$ は零でない。従って $I_6P_2$ は六次元で、その零化空間は
$T\mathbin{\lrcorner}P_2$ に次元から一致する。

$L_j=G^{-1}H_{5,3}(Y)x^j$、$\alpha_j=Y(fx^j),\ \beta_j=Y(ex^j)$ ($0\le j\le2$) とおく。
通常の多項式の等式として
\[
                L_{j+1}-xL_j=\alpha_j e-\beta_j f
\]
が成り立つ。実際 $A$ で左辺は $1,\ldots,x^4$ と直交し、$x^5$ との対合は
$x^6$ の剰余を用いて $\alpha_j$ である。また $L_j$ の五次係数は
$G(e,L_j)=\beta_j$ なので、多項式への持上げも式の通りである。
これから $K_{ij}=G(L_i,L_j)$ について
$K_{i+1,j}-K_{i,j+1}=\alpha_i\beta_j-\beta_i\alpha_j$ を得る。

反対切断を $U$ の対称二乗から $Q$ 上の対称形式への写像として読む。
$P_2\cdot P_2$ と $P_2\cdot z$ の像は六次のHankel空間内でそれぞれ
$T\mathbin{\lrcorner}P_4$ と $Y\mathbin{\lrcorner}P_2$ である。
前者は五次元で、その零化空間が $I_6$ なので、和の次元は $5+\operatorname{rank}(\alpha,\beta)$。
$z^2$ の像 $K$ は上式から階数2の場合に限りHankelでない。従ってこの場合全階数は8となる。
\end{proof}
\begin{proof}[両始点の交叉]
$Q,U,Q$ 順の三次交叉の元を、両側からそれぞれ関係 $B^*(P_5^*)$ に入る三重線型形式として読む。
$U$ を $P_2$ に制限すると、両方のHankel条件から係数は三指数の和だけに依存し、
一意な $D\in P_8^*$ になる。この制限から両側の関係係数は全て戻る
（$P_2P_3=P_5$）。$z$ で両側が一致するための残る条件は
$M_{ij}=G(L_i,L'_j)$ が対称であることであり、$L'_j$ は $D$ から同様に作る。
$D$ の二行を $\alpha',\beta'$ とすると
\[
 M_{i,j+1}-M_{i+1,j}=\beta_i\alpha'_j-\alpha_i\beta'_j.
\]
各反対角線で差を加えて $M_{ij}-M_{ji}$ を取ると、対称性は
$\beta\odot\alpha'=\alpha\odot\beta'$ と同値である。実際長さ1の差で対角三成分、
長さ2で隣の二成分、長さ3で残る一成分が順に得られる。
階数2なら独立な一次式 $\alpha,\beta$ の一意分解から
$(\alpha',\beta')$ は $(\alpha,\beta)$ のスカラー倍だけである。
従って交叉は $\dim(T\mathbin{\lrcorner}P_2)+1=4$。
階数1なら行基底を変えて $\beta=0$ とすれば $\beta'=0$ だけが条件で、交叉は六次元となる。

$U,Q,U$ 順では両端を $P_2$ に制限すると、同じ指数のずらしにより
$D\in P_7^*$ を得る。片方の端だけを $z$ にした値は
$D(x^i L_j)$ ($0\le i\le2$, $0\le j\le3$) に強制される。
これがHankelなら、端が $z$ の側の六つの関係係数も一意に定まり、両端が $z$ の値は
左右でその同じ汎関数を $L_j$ に作用させるので一致する。
従って条件は $D(x^i(\alpha_j e-\beta_j f))=0$ ($i=0,1$, $0\le j\le2$) で尽きる。
階数2なら $D$ は四次元の $eP_1+fP_1$ を零化し、解も四次元となる。

以上より四次交叉の元は、第一文字または末尾文字の自己写像 $S,N$ を $w$ に作用させたものとして
両方に表せる。例えば $w(Su,q,u',q')=w(u,q,u',Nq')$ である。
左辺の $q,q'$ 交換不変性を使えば右辺の $N$ を $q$ へ移せる。
$B$ の全射性と $G$ の非退化性から $B(Su,q)=B(u,Nq)$ となる。
さらに元の等式を使えば $B(u,Nq)$ は $B(u,q)$ に一つの線型写像を作用させたものになる。
従って $B(p,q)=pq$ に制限して $N$ は二元積の全ての一致を保つ。
$p=1,x$ により $N(x^{j+1})=xN(x^j)$ ($j=0,1,2$) となり、全ての値の次数上界3から $N$ はスカラー。
他方の向きでは役割を交換して同じ等式を得る。両四次交叉は一次元で、
復元定理が有限Frobenius性を与える。
\end{proof}
\begin{proof}[商の印の復元]
$Q$ 側の反対核は二次元で、二元三次評価のねじれ三次曲線 $C\subset\mathbb P(Q^*)$
を消す三つの二次式の中の鉛筆である。各非零式は既約である
（非退化整曲線 $C$ は二平面の和に入らない）。
従って二つの二次式は次数4の完全交叉を切り、その中で $C$ は唯一のねじれ三次曲線である。
よって $C$ とその二次Veroneseの七次元線型包であるHankel空間は内在的に定まる。
この空間で反対切断の値を割ると、$P_2\cdot U$ では零で $z^2$ だけが非零。
この階数1の二次形式の根基が元の三平面 $P_2\subset U$ を復元する。
これで商の核直線が決まり、命題\ref{prop:binary-quotient-rank-two-curve}で
商における階数一曲線と両側の二元座標が共通の射影座標変換まで定まる。
第四文字を加える自由度は前定理で既に割っている。

独立な四位置の同型については、点グラフが反転を交互に繰り返すので両偶数位置の誘導写像は
点スキーム上で一致する。点スキームの二次式制限は単射であり、二つの射影線型写像が
そこで一致すれば比例小行列式が恒等零となる。奇数位置についても同じで、
従って二位置の基底変更だけを考えればよい。逆に二元座標の変換は積、極性空間、Gram形式をすべて運ぶ。
既約なパラメータ空間の次元は $10+5=15$、一般の十次形式の射影安定群は有限なので
商の次元を得る。差分式と両階数での交叉検算は
\path{computations/verify_binary_apolar_rank_test.py}にある。
\end{proof}
